% Generated draft: requires mathematical and source review.
\documentclass[11pt]{article}
\usepackage[margin=25mm]{geometry}
\usepackage{fontspec}
\setmainfont{Latin Modern Roman}
\usepackage{amsmath,amssymb,mathtools,unicode-math}
\setmathfont{Latin Modern Math}
\usepackage{enumitem,needspace,xurl,hyperref}
\setlength{\parindent}{0pt}
\setlength{\parskip}{5pt}
\setlength{\emergencystretch}{3em}
\newcommand{\N}{\mathbb{N}}
\newcommand{\Z}{\mathbb{Z}}
\newcommand{\Q}{\mathbb{Q}}
\newcommand{\R}{\mathbb{R}}
\newcommand{\C}{\mathbb{C}}
\newcommand{\F}{\mathbb{F}}
\newcommand{\A}{\mathbb{A}}
\newcommand{\PP}{\mathbb{P}}
\newcommand{\E}{\mathbb{E}}
\newcommand{\eps}{\varepsilon}
\newcommand{\dd}{\,\mathrm{d}}
\begin{document}
\section*{UnsolvedMath ID: 701482}
\textbf{Source problem code:} catalogue entry
\subsection{Definitions and mathematical statement}
A group \(G\) is locally finite if every finitely generated subgroup of \(G\) is finite. It is simple if \(G\neq\{1\}\) and its only normal subgroups are \(\{1\}\) and \(G\). We use the standard convention that \(G\) is linear if there exist a field \(K\), an integer \(n\geq 1\), and an injective homomorphism \(G\to\operatorname{GL}_n(K)\); non-linear means that no such choices exist. The supplied record does not explicitly define linearity, so this convention requires review against the original source.

For \(g\in G\), define its centralizer by
\[
C_G(g)=\{x\in G:xg=gx\}.
\]
For a group \(H\), put \(H^{(0)}=H\) and \(H^{(j+1)}=[H^{(j)},H^{(j)}]\), where \([A,A]\) is the subgroup generated by commutators of elements of \(A\). A group \(H\) is soluble if \(H^{(d)}=\{1\}\) for some integer \(d\geq 0\). A group is almost soluble if it has a soluble subgroup of finite index.

\textbf{Problem (catalogue entry).} Does there exist a non-linear simple locally finite group \(G\) such that
\[
\forall g\in G\setminus\{1\}\quad
\exists H_g\leq C_G(g)\quad
\exists d_g\in\mathbb Z_{\geq 0}\quad
[C_G(g):H_g]<\infty
\quad\text{and}\quad
H_g^{(d_g)}=\{1\}?
\]
Neither the finite index nor the derived length is required to have a bound independent of \(g\). The supplied registry labels this problem open; that label is supplied source metadata, not an independently verified assessment.
\subsection{Short English statement}
Does a non-linear simple locally finite group exist in which the centralizer of every nonidentity element has a soluble subgroup of finite index? Both that index and the subgroup's derived length may depend on the element.
\subsection{Sources}
\begin{itemize}
\item Local UnsolvedMath registry, permanent ID \(701482\), exact source problem code \texttt{catalogue entry}, title \emph{Kourovka Notebook Problem 18.53}. The supplied record identifies upstream revision \texttt{d371ee917ac9647fbb06cd6fc529deec6f3dda82} and gives the status as open.
\item The supplied background attributes the question to \emph{The Kourovka Notebook (Unsolved Problems in Group Theory)}, 18th issue (2014), Problem 18.53, and also identifies the 21st issue edition (2026), \texttt{arXiv:1401.0300v48}. These bibliographic details are supplied attribution and have not been independently checked. The registry's collection description names E. I. Khukhro and V. D. Mazurov as editors.
\item The supplied proposer field reads \texttt{M. Kuzucuo\textbackslash u\{g}. The background repeats that truncated string and places a separate fragment \texttt{lu\}} after the discussion heading. The intended complete proposer name cannot be reliably recovered from these malformed fields without checking the source; it has not been silently repaired. The proposed year is supplied as 2014.
\end{itemize}
\end{document}
